\scanpage{172}
et si $S$ est le spectre d'un anneau de Dedekind
\[
H^1_{\unsure{\mathcal C|\mathscr C|C}_2}\xrightarrow{\sim}H^1_{\mathcal C_1}
\]
[d'où dans ce cas $H^1_{\mathcal C_3}\xrightarrow{\sim}H^1_{\mathcal C_2}\simeq H^1_{\mathcal C_1}$]. Si $S$ est le spectre d'un anneau \deleted{local complet}\inserted{de valuation} discrète à corps résiduel alg. clos, \inserted{$V$} on aura de $\overline m$
\[
H^1_{\mathcal C_i}\unsure{b|6}\simeq H^1_{\deleted{\unsure{\mathcal C_i|4}}}(\widehat V/V,-)\qquad\text{pour }i=1,2,3
\]
Si $S$ est le spectre d'un anneau local \uline{complet} $V_{\unsure{\phi|0}}$, alors
\[
\deleted{H^1_{\unsure{\mathcal C_4|4}}(\unsure{V|k},\unsure{-|\,})\simeq}
\]
\[
H^1_{\mathcal C_5}\simeq H^1_{\mathcal C_4}\simeq H^1_{\mathcal C_3}\quad(\simeq H^1_{\mathcal C_2})
\]
et on aura de plus
\[
H^1_{\mathcal C_i}(V,-)\simeq H^1_{\mathcal C_i}(\underset{\substack{\downarrow\\\text{corps}\\\text{résiduel}}}{k},-)\qquad(i=2,3,4,5)
\]
% archive sheet number: 171
